% deck.tex — an explorable deck in beamer. Build with: sh build.sh
% (lualatex; the theme loads the fonts TeX Live ships). Every frame below
% demonstrates one layout from theme/explorable.sty. Replace the content, keep
% the patterns: one idea per frame, title as an assertion, a \note on every
% frame that tells the presenter what to do and say.
\documentclass[aspectratio=169, 10pt]{beamer}
\usepackage{theme/explorable}
% data/numbers.tex — the single home of every number the deck uses.
%
% Scalars:  \defnum{path}{value}        read with \D{path} in prose, \Dv{path} in charts
% Tables:   \pgfplotstableread[col sep=comma]{data/name.csv}\name
%           read with \Dt{\name}{row}{col} in prose, \addplot table {\name} in charts
%
% Say in a comment where each number comes from (table, column, row). Prose
% numbers that are not in a chart are transcribed from the same tables, never
% from another slide. When a source number changes, change it here and grep
% deck.tex for the old value.
%
% If an explorable already exists for this topic, transcribe its data object
% here (and port its model functions to pgfmath) so the two never disagree.

% ---- stack-reveal: two cases, each with two components of a total (units per year)
% Source: <table and column>. Slopes are the source's; levels are stylised around sample means.
\defnum{reveal/unit}{\% per year}
\defnum{reveal/part/1}{first component}
\defnum{reveal/part/2}{second component}
\defnum{reveal/a/name}{Case A}
\defnum{reveal/a/1}{2.8}
\defnum{reveal/a/2}{5.8}
\defnum{reveal/b/name}{Case B}
\defnum{reveal/b/1}{4.1}
\defnum{reveal/b/2}{4.5}
\defnum{reveal/gap}{1.4}            % a/1 − b/1 rounded: the number the prose quotes

% ---- line-handle: y = a + b·(x − x0) over [xmin, xmax]; the handle stops at three x values
% Source: <table, specification>. The slope b is the estimate; a is an illustrative level.
\defnum{line/a}{6.0}
\defnum{line/b}{0.45}
\defnum{line/x0}{10}
\defnum{line/xmin}{2}
\defnum{line/xmax}{18}
\defnum{line/ymin}{0}
\defnum{line/ymax}{12}
\defnum{line/stops/1}{5}
\defnum{line/stops/2}{10}
\defnum{line/stops/3}{15}
\defnum{line/xlabel}{position (units)}
\defnum{line/ylabel}{outcome (units)}

% ---- ladder: an estimate as controls are added (cumulative down the rows)
% Columns: spec label, beta, se, sig (1 if the source's table marks it at 5 %, else 0).
% Significance comes from the table's stars, not from rounded standard errors.
\pgfplotstableread[col sep=comma]{data/ladder.csv}\ladder
\defnum{ladder/xlabel}{coefficient on X (units)}
\defnum{ladder/span}{3 to 1}       % the span between the two cases, quoted on the stats slide

% ---- stylised dots for the line-handle stage (seeded draw around the line; see data/README)
\pgfplotstableread[col sep=comma]{data/line-dots.csv}\linedots


% Topic aliases for the four colour tokens, one line per token, in the topic's
% words. Charts and prose use the aliases; the generic names stay.
%   \colorlet{gain}{primary}  \colorlet{yield}{contrast}  \colorlet{anna}{actora}  \colorlet{bo}{actorb}

\title{Deck title as a claim}
\subtitle{Venue, month year}
\author{Author One ~·~ Author Two}
\institute{Institution ~·~ Institution}
\date{}

% \usepackage[style=authoryear]{biblatex} \addbibresource{references.bib}   % if the deck cites

\begin{document}

\begin{frame}[plain]
\titlepage
\note{Slide 1 is the title: screenshot numbers, plan numbers and note numbers all start here.}
\end{frame}

\begin{frame}{The question the room wants answered before you answer it}
\begin{prose}
One or two sentences that make the room want the answer.
\end{prose}
\note{The hook. Ask the room to commit: a show of hands, one sentence. Then advance to the payoff.}
\end{frame}

\begin{frame}{Two cases, one reveal}
\begin{beside}[c]
\begin{prose}
\textbf{Case A earned more on the first component.} About \num{\primaryt{\D{reveal/gap} points}} a year.

\textbf{Case B earned more on the second.} The two components offset, leaving the totals \contrastt{about equal}.
\end{prose}
\nextcol
\begin{stage}
% charts/stack-reveal.tex — two cases, two components each. The second
% component builds in on the second overlay: the payoff after a bet.
% The frame needs no overlay spec of its own; visible on=<2-> creates step 2.
% Numbers: data/numbers.tex, reveal/*.
\pgfmathsetmacro{\revealTotA}{\Dv{reveal/a/1}+\Dv{reveal/a/2}}
\pgfmathsetmacro{\revealTotB}{\Dv{reveal/b/1}+\Dv{reveal/b/2}}
\begin{tikzpicture}
\begin{axis}[deckaxis, width=72mm, height=44mm,
  ybar stacked, bar width=16mm, enlarge x limits=false,
  xmin=0.4, xmax=2.6, xtick={1,2}, xticklabels={\Dv{reveal/a/name}, \Dv{reveal/b/name}},
  xticklabel style={font=\sffamily\bfseries\fontsize{7.5}{9}\selectfont, color=ink},
  ymin=0, ymax=11.5, ytick={0,4,8}, ylabel={\Dv{reveal/unit}},
  x axis line style={draw=none}, xtick style={draw=none},
]
% first component, solid
\addplot[fill=actora, draw=none] coordinates {(1,\Dv{reveal/a/1}) (2,0)};
\addplot[fill=actorb, draw=none] coordinates {(1,0) (2,\Dv{reveal/b/1})};
% second component, hatched, builds in on step 2
\addplot[pattern=north east lines, pattern color=actora, draw=actora, line width=0.6pt, visible on=<2->] coordinates {(1,\Dv{reveal/a/2}) (2,0)};
\addplot[pattern=north east lines, pattern color=actorb, draw=actorb, line width=0.6pt, visible on=<2->] coordinates {(1,0) (2,\Dv{reveal/b/2})};
% labels: the first component on step 1, the totals on step 2
\node[datalabel, anchor=south, text=actora, visible on=<1>] at (axis cs:1,\Dv{reveal/a/1}) {\Dv{reveal/a/1}};
\node[datalabel, anchor=south, text=actorb, visible on=<1>] at (axis cs:2,\Dv{reveal/b/1}) {\Dv{reveal/b/1}};
\node[datalabel, anchor=south, text=actora, visible on=<2->] at (axis cs:1,\revealTotA) {\pgfmathprintnumber[fixed,precision=1]{\revealTotA}};
\node[datalabel, anchor=south, text=actorb, visible on=<2->] at (axis cs:2,\revealTotB) {\pgfmathprintnumber[fixed,precision=1]{\revealTotB}};
\end{axis}
\end{tikzpicture}
\legend{\swatch{ink2}~\Dv{reveal/part/1}\qquad{\tikz\fill[pattern=north east lines, pattern color=ink2] (0,0) rectangle (1.6ex,1.6ex);}~\Dv{reveal/part/2}}
\stagecap{Slopes are the source's estimates; levels are stylised around sample averages (\Dv{reveal/unit}).}

\end{stage}
\end{beside}
\note{Step 1 shows the first component only. Ask for the bet, then advance: the second component builds in. Let it land before speaking.}
\end{frame}

\begin{frame}{What the literature measures, and what this adds}
\begin{claims}
\claim{First strand.}{What it did and what it could not see.}
\claim{Second strand.}{Where it stopped.}
\claim{This talk.}{The object we observe that they could not.}
\end{claims}
\note{One slide, three strands at most. Say why each stopped where it did.}
\end{frame}

\divider[Section 1]{The first section as a claim}

\begin{frame}[label=gradient]{The gradient in one picture}
\begin{wideleft}[c]
\fig{figs/example.pdf}
\nextcol
\begin{prose}
One sentence about what the figure shows.

\begin{stats}[1]
\stat{\Dt[3]{\ladder}{0}{beta}}{unit, and the magnitude between two points the room can picture. The number is read from data/numbers.tex.}
\end{stats}
\end{prose}
\source{What is plotted, the sample, the source table or figure.}
\end{wideleft}
\note{Replace figs/example.pdf with a source figure: sh figs/convert\_figs.sh <source-dir> <Name>. The side room on robustness opens from here.}
\end{frame}

\begin{frame}<1-3>{A line the presenter can walk along}
\begin{beside}[c]
\begin{prose}
Step through the three stops. The card says what a unit at that position earns.

\source{Dots are stylised around the estimated line. The slope is the source's estimate.}
\end{prose}
\nextcol
\begin{stage}
% charts/line-handle.tex — a fitted line with a handle the presenter walks
% along. Three stops, one per overlay step: \stagesteps{3} here, <1-3> on the
% frame. Step 1 is where the picture is most interesting; in the handout the
% handle sits at the last stop. Numbers: data/numbers.tex, line/*.
\stagesteps{3}
\pgfmathsetmacro{\handleX}{\Dv{line/stops/\step}}
\pgfmathsetmacro{\handleY}{\Dv{line/a}+\Dv{line/b}*(\handleX-\Dv{line/x0})}
\begin{tikzpicture}
\begin{axis}[deckaxis, width=\chartwidth, height=\chartheight,
  xmin=\Dv{line/xmin}, xmax=\Dv{line/xmax}, ymin=\Dv{line/ymin}, ymax=\Dv{line/ymax},
  xlabel={\Dv{line/xlabel}}, ylabel={\Dv{line/ylabel}},
]
\addplot[only marks, mark=*, mark size=1.1pt, mark options={fill=line, draw=none}] table[x=x, y=y] {\linedots};
\addplot[primary, domain=\Dv{line/xmin}:\Dv{line/xmax}, samples=2] {\Dv{line/a}+\Dv{line/b}*(x-\Dv{line/x0})};
\draw[line, dashed, line width=0.5pt] (axis cs:\handleX,\Dv{line/ymin}) -- (axis cs:\handleX,\handleY);
\node[handle] (h) at (axis cs:\handleX,\handleY) {};
\node[card, above=2.5mm of h] {at \pgfmathprintnumber[fixed,precision=0]{\handleX}: \textbf{\pgfmathprintnumber[fixed,precision=1]{\handleY}} units};
\end{axis}
\end{tikzpicture}
\stagecap{Dots are stylised around the estimated line. The slope (\D[2]{line/b} per unit) is the source's estimate; the level is illustrative.}

\end{stage}
\end{beside}
\note{Three steps, one per stop. Stop at the two points named on the previous slide and read the card aloud. If the projector drops a step, say the three numbers.}
\end{frame}

\begin{frame}{Adding controls removes most of the slope}
\begin{beside}[c]
\begin{prose}
Controls are cumulative down the rows. The last row is statistically indistinguishable from zero.

\source{Coefficient on X with 95 percent intervals. Specification details. \goto{robust}{robustness}}
\end{prose}
\nextcol
\begin{stage}
% charts/ladder.tex — a coefficient as controls are added, one row per
% specification, 95 percent intervals. A row the source's table does not mark
% as significant (sig = 0 in data/ladder.csv) is drawn grey. The two addplots
% filter the same table by sig, so rows keep their order and colour their bars.
\pgfplotstablegetrowsof{\ladder}\pgfmathtruncatemacro{\ladderRows}{\pgfplotsretval}
\pgfmathtruncatemacro{\ladderTop}{\ladderRows-1}
\begin{tikzpicture}
\begin{axis}[deckaxis, width=\chartwidth, height=\chartheight,
  xlabel={\Dv{ladder/xlabel}},
  ymin=-0.6, ymax=\ladderTop+0.6, y dir=reverse,
  ytick={0,...,\ladderTop}, yticklabels from table={\ladder}{spec},
  yticklabel style={font=\sffamily\fontsize{7}{8.5}\selectfont, color=ink},
  y axis line style={draw=none}, ytick style={draw=none},
  ymajorgrids=false, xmajorgrids,
  error bars/x dir=both, error bars/x explicit, error bars/error mark=none,
  unbounded coords=discard, filter discard warning=false,
]
\draw[line, line width=0.6pt] (axis cs:0,-0.6) -- (axis cs:0,\ladderTop+0.6);
\addplot[only marks, mark=*, mark size=2.4pt, primary, error bars/error bar style={line width=0.9pt, primary}]
  table[x=beta, y expr=\coordindex, x error expr=1.96*\thisrow{se}, restrict expr to domain={\thisrow{sig}}{1:1}] {\ladder};
\addplot[only marks, mark=*, mark size=2.4pt, ink3, error bars/error bar style={line width=0.9pt, ink3}]
  table[x=beta, y expr=\coordindex, x error expr=1.96*\thisrow{se}, restrict expr to domain={\thisrow{sig}}{0:0}] {\ladder};
\end{axis}
\end{tikzpicture}
\stagecap{Point estimate with 95 percent interval; controls are cumulative down the rows. Grey: not distinguishable from zero in the source's table.}

\end{stage}
\end{beside}
\note{Read the rows top to bottom. The point is the last row, not the first.}
\end{frame}

\begin{frame}{Three numbers the room should leave with}
\begin{stats}[3]
\stat{\D{reveal/gap} pp}{a year, between the two cases}
\stat[ink2]{\Dv{ladder/span}}{the span those two cases sit apart}
\stat[contrast]{$\approx 0$}{the total gradient, statistically}
\end{stats}
\note{A stat-card slide breaks a run of prose-beside-figure layouts. Three numbers at most, each with the magnitude a listener can feel.}
\end{frame}

\assertion{The one sentence: the return arrives in different forms, and where you can buy decides the form}

\begin{frame}{Two readings of the same fact}
\begin{cols}
\colhead{Reading \actorat{one}}
\begin{prose}
Two sentences on what it predicts.
\end{prose}
\nextcol
\colhead{Reading \actorbt{two}}
\begin{prose}
Two sentences on what it predicts instead.
\end{prose}
\end{cols}
\vspace{4mm}
\q{The test that separates them: what we look at next.}
\note{Name the two readings, then the one observation that tells them apart.}
\end{frame}

\begin{frame}{Three things to take away}
\begin{claims}
\claim{}{The first finding, with its magnitude.}
\claim{}{The second, and why it reverses the obvious reading.}
\claim{}{What it implies, in one sentence.}
\end{claims}
\note{End on the room's question, not yours.}
\end{frame}

\begin{frame}{Closing claim, the sentence they should leave with}
\begin{prose}
\textbf{Thank you.}
\end{prose}
\source{Name}\source{Affiliation}\source{email}
\note{Leave this up during questions; the rooms are one link away.}
\end{frame}

% ---------------------------------------------------------------- appendix
% Frames after \appendix are numbered A1, A2, … and leave the progress bar,
% so the count ends on the closing slide.
\appendix
\divider{Appendix}

\begin{frame}[label=robust]{A side room: the robustness a questioner will ask about}
\figpair{figs/example.pdf}{figs/example.pdf}
\figcap{\textbf{Left:} what. \textbf{Right:} what. Source.}
\back{gradient}
\note{Opened from the ladder slide. Say what changes and what does not, then go back.}
\end{frame}

% \begin{frame}[allowframebreaks]{References}\printbibliography\end{frame}

\end{document}
